% !TeX root = ../main.tex

\xchapter{总结与展望}{Conclusions and Outlook}
\label{chap:conclusions_outlook}

\xsection{主要研究工作与结论}{Main Research Work and Conclusions}

快堆MOX燃料元件在高功率密度、高温度梯度和强中子辐照条件下服役，其温度、组分、微观结构、裂变气体行为以及力学响应之间存在显著的双向耦合关系。针对传统燃料性能分析中材料模型分散、微观演化过程考虑不足、多物理场耦合求解困难以及高保真模型在线计算代价较高等问题，本文以快堆MOX燃料元件为研究对象，围绕“材料模型构建—微观结构与组分演化—非线性多物理场耦合实现—燃料棒级验证—完整场快速预测”开展了系统研究。主要研究工作与结论如下。

\begin{enumerate}[label={（\arabic*）},leftmargin=2.8em]
	\item 建立了适用于氧化物燃料元件分析的材料模型体系。系统梳理并实现了MOX燃料热导率、比热容、热膨胀、弹性、蠕变、密实化、肿胀和重定位等模型，同时整理了SS316、D9和HT9等典型快堆包壳材料的热物性、塑性、蠕变、辐照肿胀和损伤模型。在燃料--包壳间隙换热方面，建立了由混合气体传导、固体接触传导和热辐射组成的统一模型，并明确了温度、压力、气体组分、间隙宽度和接触压力等状态变量之间的传递关系。该部分工作为后续微观演化、裂变气体和热--力耦合分析提供了统一的物性基础。

	\item 建立了考虑气孔迁移、锕系元素重分布和氧金属比演化的强耦合模型。针对高温度梯度下气孔向燃料中心迁移并引起组织重构的问题，构建了孔隙输运方程，比较了经验关系式与Rand--Markin多组分热力学模型对气孔迁移速率的描述，并通过稳定化离散提高了对流占优条件下的数值稳定性。针对钚重分布，建立了同时考虑化学扩散、热扩散和气孔迁移增强效应的组分输运模型，能够再现钚向高温区域富集的主要趋势，并进一步扩展至镅元素迁移分析。针对O/M演化，建立了热扩散与化学扩散耦合的氧迁移模型。计算结果表明，初始化学计量状态会显著影响氧的径向迁移方向和幅度，而局部O/M变化又会通过蒸气压、热导率、熔点及迁移驱动力反馈影响气孔迁移和锕系元素重分布。

	\item 完成了非线性力学、裂变气体行为和间隙导热模型的统一数值实现。针对燃料和包壳材料中塑性、蠕变及多种非弹性应变共同作用的问题，建立了热弹塑性控制方程的有限元离散形式，并采用径向返回算法在材料积分点处隐式更新应力、塑性应变和硬化变量。裂变气体方面，分别实现了适用于工程计算的ForMas模型和细观机理程序SCIANTIX的耦合接口，实现了温度、裂变率、燃耗、应力及气体状态变量的双向传递。进一步将裂变气体释放、气腔压力、混合气体组分和间隙换热系数纳入统一时间推进过程。单元级、子模型级和耦合接口级验证结果表明，所建立的数值算法能够保持主要物理趋势和状态变量演化的一致性，为燃料棒长期热--力耦合计算提供了稳定的计算基础。

	\item 建立了面向MOX燃料棒的验证流程。基于公开基准数据和补充实验数据，构建了燃料中心温度与裂变气体释放的棒级计算模型，并分别采用ForMas和SCIANTIX描述裂变气体行为。已完成的中心温度算例能够总体再现实验数据随功率和燃耗变化的主要趋势，说明所建立模型能够描述MOX燃料棒的稳态热响应。高燃耗阶段局部偏差增大，表明燃料热导率退化、初始和演化间隙、裂变气体释放以及接触换热仍是影响预测精度的关键因素。相关分析同时说明，温度、裂变气体释放和间隙传热之间应作为闭环反馈过程处理，而不宜采用彼此独立的单向计算方式。

	\item 构建了基于POD--LSTM的燃料性能完整场快速预测方法。以经过验证的二维RZ燃料棒模型生成高保真场快照，通过网格敏感性分析选取包含5717个节点和5000个单元的统一网格。采用POD将温度、径向位移和轴向位移共17151维节点输出压缩为28维低阶系数，并首先以插值和POD--GPR模型建立固定功率倍率下的预测基准。在此基础上，构建了包含40条训练历史、10条验证历史和10条独立盲测历史的变功率数据集，采用时间、瞬时功率、功率变化率和累积功率积分作为序列特征，利用LSTM预测POD系数增量并恢复完整二维场。独立盲测中，温度、径向位移和轴向位移的平均相对$L_2$误差分别为1.256\%、1.528\%和2.549\%。单条109步历史的系数预测与完整场重构平均耗时为0.0743~s，独立在线示例的总计算时间为0.06203~s，表明该方法能够在训练数据包络内快速获得全寿期温度场和位移场。当前尚未补充对应高保真工况的统一壁钟时间，因此本文不对加速比作定量结论。
\end{enumerate}

\xsection{主要创新点}{Main Innovations}

本文的主要创新点概括如下。

\begin{enumerate}[label={（\arabic*）},leftmargin=2.8em]
	\item 建立了快堆MOX燃料微观结构重构、多组分迁移与氧化学状态演化的统一耦合模型。区别于将气孔迁移、钚重分布和O/M演化分别处理的传统建模方式，本文以温度场、氧分压、蒸气压和局部材料状态为共同耦合变量，在统一有限元框架中描述孔隙迁移、中心孔形成、钚及镅再分布和氧迁移过程。模型同时考虑固相扩散、热扩散及气孔迁移增强效应，并将O/M对热物性和迁移驱动力的反馈纳入计算，从而形成“微观结构—组分分布—热化学状态”相互作用的闭环描述，为分析高温度梯度和高燃耗条件下MOX燃料的组织与组分演化提供了更完整的模型基础。

	\item 构建了非线性力学、细观裂变气体行为与间隙传热相耦合的隐式数值实现框架。本文将基于$J_2$塑性的径向返回算法用于燃料和包壳多非弹性本构的局部隐式积分，并在统一时间推进过程中耦合ForMas工程模型、SCIANTIX细观机理模型及间隙换热模型。通过外部机理程序状态变量的持续传递与回写，使晶内和晶界气体演化、气体释放、气腔压力、混合气体导热、间隙闭合及力学变形能够相互反馈。所建立的“局部本构积分—细观状态更新—全局热力平衡”分层计算方式，提高了复杂燃料性能模型的模块化程度，并为长期辐照条件下多物理场强耦合计算提供了稳定的实现途径。

	\item 提出了变功率历史驱动的POD--LSTM燃料性能完整场快速预测方法。该方法不是仅对最高温度或最大位移等标量进行回归，而是首先利用POD提取温度场、径向位移场和轴向位移场的低维空间模态，再利用具有历史记忆能力的LSTM预测低维系数随功率历史的演化，并在任意输出时刻恢复完整二维节点场。模型采用系数增量预测、可训练循环初始状态以及按完整功率历史划分训练、验证和盲测数据的策略，降低了初始状态误差和时间点随机划分引起的数据泄漏风险。独立盲测和在线示例表明，该方法能够在训练包络内以百毫秒以内的计算代价获得全寿期完整热力场，为高保真燃料性能模型由离线计算向快速运行历史分析和在线状态评估扩展提供了技术途径。
\end{enumerate}

\xsection{研究展望}{Research Outlook}

尽管本文围绕快堆MOX燃料材料模型、微观演化、多物理场耦合和快速预测开展了系统研究，但受实验数据、参数覆盖范围和计算成本限制，仍有若干问题需要进一步研究。

\begin{enumerate}[label={（\arabic*）},leftmargin=2.8em]
	\item 进一步完善微观结构与多组分迁移模型。当前气孔迁移、钚及镅重分布和O/M演化模型中的部分扩散系数、气相迁移参数及热力学参数仍具有较大不确定性。后续可结合辐照后径向组分测量、孔隙和晶粒组织表征以及中心孔演化数据开展联合标定，并引入参数不确定性传播和敏感性分析，定量评估微观模型对中心温度、熔化裕度和局部功率分布的影响。同时，可进一步考虑晶粒长大、高燃耗结构、裂纹形成及裂纹与孔隙迁移之间的相互作用。

	\item 扩展燃料--包壳强耦合与失效分析能力。现有力学框架主要建立在小变形与连续介质假设下，后续可引入有限变形、芯块开裂、碎块重定位、摩擦接触和更精细的PCMI模型，并结合包壳辐照硬化、损伤、腐蚀及失效判据研究稳态和瞬态工况下的安全裕度。对于裂变气体模型，还需要进一步提高SCIANTIX在高燃耗、非化学计量MOX及复杂应力状态下的适用性，并加强裂变气体肿胀、燃料变形和间隙传热之间的一致性处理。

	\item 完善燃料棒级验证与不确定性评价。后续应继续补充快堆MOX燃料及不同包壳材料的中心温度、裂变气体释放、包壳应变、内压和辐照后检查数据，形成覆盖稳态辐照、功率提升和典型瞬态工况的分层验证矩阵。在比较计算值与实验值的同时，应系统考虑功率历史、初始间隙、晶粒尺寸、充气压力、轴向功率分布及实验测量误差，建立模型参数、输入条件和实验数据共同作用下的不确定性量化方法。

	\item 扩展快速预测模型的参数空间和可信性边界。当前POD--LSTM模型限定于固定几何、固定材料模型和固定冷却剂边界条件下的变功率历史预测。后续可将冷却剂温度和压力、填充气体压力、燃料密度、初始间隙及几何参数纳入输入空间，并通过参考域映射或网格无关表示处理不同几何之间的场数据。还可引入物理约束损失、守恒关系、代理模型不确定性估计和超出训练包络检测机制，提高模型在复杂运行历史下的可靠性。

	\item 探索高保真模型与运行数据融合的在线应用。可在现有非侵入式降阶框架基础上进一步结合传感器数据、数据同化和在线参数更新方法，使代理模型能够根据实测功率、冷却剂条件及可观测响应持续修正内部状态。通过建立高保真模型、快速代理模型和运行数据之间的闭环更新机制，可为燃料元件数字孪生、运行状态评估和异常工况快速分析提供进一步支撑。
\end{enumerate}
